\documentclass{article}
\usepackage{graphicx, physics, dsfont, braket}
\usepackage{fancyhdr}
\usepackage{hyperref}
\usepackage{ragged2e}
\usepackage{amsmath, amsthm}
\usepackage[margin=1in]{geometry}
\usepackage{setspace}
\usepackage{tikz}
\usetikzlibrary{positioning}
\pagestyle{fancy}
\lhead{Zoeb Izzi}
\rhead{\today}
\lfoot{Chapter who cares}
\rfoot{}
\begin{document}
\thispagestyle{fancy}
\begin{center}\LARGE{Electrodynamics \\ Notes}\end{center}
\section{Recap/Intro}
\section{Today's stuff}
Quantum Harmonic Oscillator.
\\\\
We start given that some $\hat H\phi = E_n\phi$. From this, we know that:
\[\hbar \omega (\hat a _+ \hat a _- + \frac 1 2)\phi = E_n \phi\]
By the raising operator, 
\[\hbar \omega(a_+ a_- a_+ + \frac 1 2 a_+) \phi = E_n(a_+ \phi)\]
\[\hbar \omega a_+ (a_- a_+ + \frac 1 2)\phi = a_+ \left[ \hbar \omega (a_- a_+ + \frac 1 2)\right]\phi\]
\[\hbar \omega (a_-a_+ - \frac 1 2) = H\]
\[a_+(\hat H + \hbar \omega)\phi = a_+ (E+\hbar \omega) \phi = (E+\hbar \omega)a_+ \phi\]
Thus, 
\[H(a_+\phi) = (E+\hbar \omega)(a_+\phi)\]
from which we can find that:
\[E_n = \hbar \omega\left(\frac 1 2 + n\right)\]
Further, we MUST REMEMBER THAT:
\[\hat a_+ \phi_n = \sqrt{n+1} \phi_{n+1}\]
\[\hat a_- \phi_n = \sqrt{n} \phi_{n-1}\]
\[\hat a_- \hat a_+ \phi_n = n \phi _n\]
What do overlap integrals calculate? Know how to calculate them and find out what you are trying to match.

\end{document}